% One-time imported draft. Editable native source; no Markdown regeneration.
\section{A bit pattern can name different
integers}\label{004-bit-inspector-p00-a-bit-pattern-can-name-different-integers}

The sprite chapter treated a byte as eight independent pixel decisions.
This chapter treats a word as an integer while keeping its underlying
bits visible. There are 2\^{}32 possible word patterns. An unsigned
interpretation assigns them the values zero through 4294967295. A signed
two\textquotesingle s-complement interpretation assigns the same
patterns the values \ensuremath{-}2147483648 through 2147483647.

The pattern with every bit set illustrates the difference. Its unsigned
value is the sum of all 32 positive binary place values, 2\^{}32 \ensuremath{-} 1.
Under the signed interpretation, the highest bit has weight \ensuremath{-}2\^{}31
instead of positive 2\^{}31. Together with the remaining set bits, that
produces \ensuremath{-}1.

The input is a word in memory, so the test can substitute any of its
possible patterns without relying on a console parser\textquotesingle s
accepted decimal range. The report shows four representations of that
original word. Separate full-width results show what happens when it is
loaded or used in arithmetic.
